% GENERATED FILE. Edit canonical lessons, then run npm run generate:books.
% canonical-source-hash: fnv1a64:f568945be5ffe3aa
% canonical-lessons: LA-C63-bracchium, LA-C63-digitus, LA-C63-dens, LA-C63-lingua, LA-C63-nasus

\chapter{Arm, Finger, Tooth}
\label{ch:la-arm-finger-tooth}
\hlchaptermodality{63}

\begin{chapteropening}
\textbf{By the end of this chapter:} \emph{I can name an arm, a finger, a tooth, the tongue and the nose.}

Complete the last lesson of chapter 63: I can name an arm, a finger, a tooth, the tongue and the nose.
\end{chapteropening}

\section[bracchium, bracchiī]{bracchium, bracchiī — an arm}

\label{lesson:LA-C63-bracchium}



\begin{quote}
Before the new one: say the Latin for before, then the Latin for after.
\end{quote}

\subsection*{You'll want to know: bracchium, bracchiī}
\textbf{bracchium, bracchiī} — \textquotedblleft{}an arm\textquotedblright{}.

\begin{cousinweb}
\textbf{bracchium} is usually taken as a loan from Greek \emph{brakhiōn}, \textquotedblleft{}arm\textquotedblright{} (EtymDB 2.1). Lewis \& Short only call the two words kindred. English \emph{bracelet} comes from \emph{bracchium} through French.
\end{cousinweb}

\begin{grammarlens}[title={What you've built}]
One more word for this chapter.
\end{grammarlens}

\subsection*{Your turn}
\begin{itemize}
\raggedright
  \item \emph{Say it:} \emph{bracchium}
  \item \emph{Say it:} \emph{bracchium}, once more
  \item \emph{Say it:} say \emph{bracchium}
  \item \emph{Recall:} say the Latin for before, then the Latin for after, then say \emph{bracchium} again
\end{itemize}

\begin{tcolorbox}[breakable,colback=teal!4,colframe=teal!35!black,title={Before you move on}]
What does bracchium, bracchiī mean? (\textquotedblleft{}An arm\textquotedblright{}.) Say it once more.
\end{tcolorbox}
\section[digitus, digitī]{digitus, digitī — a finger}

\label{lesson:LA-C63-digitus}



\begin{quote}
Before the new one: say the Latin for after, then the Latin for an arm.
\end{quote}

\subsection*{You'll want to know: digitus, digitī}
\textbf{digitus, digitī} — \textquotedblleft{}a finger\textquotedblright{}. English \emph{digit} keeps it: we first counted on our fingers.

\begin{cousinweb}
\textbf{digitus} has no certain origin. It is often linked to Proto-Indo-European \emph{*deik-}, \textquotedblleft{}to show, to point out\textquotedblright{} (EtymDB 2.1): the finger as the pointer.
\end{cousinweb}

\begin{grammarlens}[title={What you've built}]
One more word for this chapter.
\end{grammarlens}

\subsection*{Your turn}
\begin{itemize}
\raggedright
  \item \emph{Say it:} \emph{digitus}
  \item \emph{Say it:} \emph{digitus}, once more
  \item \emph{Say it:} say \emph{digitus}
  \item \emph{Recall:} say the Latin for after, then the Latin for an arm, then say \emph{digitus} again
\end{itemize}

\begin{tcolorbox}[breakable,colback=teal!4,colframe=teal!35!black,title={Before you move on}]
What does digitus, digitī mean? (\textquotedblleft{}A finger\textquotedblright{}.) Say it once more.
\end{tcolorbox}
\section[dēns, dentis]{dēns, dentis — a tooth}

\label{lesson:LA-C63-dens}



\begin{quote}
Before the new one: say the Latin for an arm, then the Latin for a finger.
\end{quote}

\subsection*{You'll want to know: dēns, dentis}
\textbf{dēns, dentis} — \textquotedblleft{}a tooth\textquotedblright{}. A \emph{dentist} looks after them.

\begin{cousinweb}
\textbf{dēns} is inherited, usually traced to Proto-Indo-European \emph{*h3dent-}, \textquotedblleft{}tooth\textquotedblright{} (EtymDB 2.1). English \emph{tooth} is a distant cousin (Lewis \& Short), and English \emph{dental} comes from \emph{dēns} itself.
\end{cousinweb}

\begin{grammarlens}[title={What you've built}]
One more word for this chapter.
\end{grammarlens}

\subsection*{Your turn}
\begin{itemize}
\raggedright
  \item \emph{Say it:} \emph{dēns}
  \item \emph{Say it:} \emph{dēns}, once more
  \item \emph{Say it:} say \emph{dēns}
  \item \emph{Recall:} say the Latin for an arm, then the Latin for a finger, then say \emph{dēns} again
\end{itemize}

\begin{tcolorbox}[breakable,colback=teal!4,colframe=teal!35!black,title={Before you move on}]
What does dēns, dentis mean? (\textquotedblleft{}A tooth\textquotedblright{}.) Say it once more.
\end{tcolorbox}
\section[lingua, linguae]{lingua, linguae — the tongue; a language}

\label{lesson:LA-C63-lingua}



\begin{quote}
Before the new one: say the Latin for a finger, then the Latin for a tooth.
\end{quote}

\subsection*{You'll want to know: lingua, linguae}
\textbf{lingua, linguae} — \textquotedblleft{}the tongue; a language\textquotedblright{}. Both meanings, as in English \emph{tongue}.

\begin{cousinweb}
\textbf{lingua} is inherited. Early Latin had \emph{dingua}, \textquotedblleft{}tongue\textquotedblright{} (Lewis \& Short), and the word is usually traced to Proto-Indo-European \emph{*dnghweh2s}, \textquotedblleft{}tongue\textquotedblright{} (EtymDB 2.1); English \emph{tongue} is a cousin. The \emph{l} is often put down to the pull of \emph{lingō}, \textquotedblleft{}to lick\textquotedblright{}.
\end{cousinweb}

\begin{grammarlens}[title={What you've built}]
One more word for this chapter.
\end{grammarlens}

\subsection*{Your turn}
\begin{itemize}
\raggedright
  \item \emph{Say it:} \emph{lingua}
  \item \emph{Say it:} \emph{lingua}, once more
  \item \emph{Say it:} say \emph{lingua}
  \item \emph{Recall:} say the Latin for a finger, then the Latin for a tooth, then say \emph{lingua} again
\end{itemize}

\begin{tcolorbox}[breakable,colback=teal!4,colframe=teal!35!black,title={Before you move on}]
What does lingua, linguae mean? (\textquotedblleft{}The tongue; a language\textquotedblright{}.) Say it once more.
\end{tcolorbox}
\section[nāsus, nāsī]{nāsus, nāsī — the nose}

\label{lesson:LA-C63-nasus}



\begin{quote}
Before the new one: say the Latin for a tooth, then the Latin for the tongue; a language.
\end{quote}

\subsection*{You'll want to know: nāsus, nāsī}
\textbf{nāsus, nāsī} — \textquotedblleft{}the nose\textquotedblright{}. English \emph{nasal}.

\begin{cousinweb}
\textbf{nāsus} is inherited, usually traced to Proto-Indo-European \emph{*neh2s-}, \textquotedblleft{}nose\textquotedblright{} (EtymDB 2.1). English \emph{nose} is a cousin from the same root.
\end{cousinweb}

\begin{grammarlens}[title={What you've built}]
One more word for this chapter.
\end{grammarlens}

\subsection*{Your turn}
\begin{itemize}
\raggedright
  \item \emph{Say it:} \emph{nāsus}
  \item \emph{Say it:} \emph{nāsus}, once more
  \item \emph{Say it:} say \emph{nāsus}
  \item \emph{Recall:} say the Latin for a tooth, then the Latin for the tongue; a language, then say \emph{nāsus} again
\end{itemize}

\begin{tcolorbox}[breakable,colback=teal!4,colframe=teal!35!black,title={Before you move on}]
What does nāsus, nāsī mean? (\textquotedblleft{}The nose\textquotedblright{}.) Say it once more.
\end{tcolorbox}
